\section{Growing finite models and approximate commutation}
\label{sec:refinement}

Coupled contraction solves a fixed finite-dimensional problem. Many
mathematical constructions require another step: a solver on a small
model must agree with a solver on a larger model after restriction.
Exact commutation is often false because of boundary conditions,
truncation, or quadrature. This section derives the needed agreement
from a quantitative commutation defect.

\subsection{Finite-level data}

For each level \(n\ge1\), let \(D_n\) be an explicitly coded domain of
finite rational data, and let \(d_n\) be a rational-valued metric on it.
A pseudometric also suffices if equality of observations is understood
modulo zero distance. The domain predicate, metric, updates, restrictions, returned points
\(y_n\), and numerical bounds \(q_n,L_n,\eta_n,r_n\) are fixed PR families. Their algebraic properties are
proved locally or retained as explicit hypotheses.

Suppose
\[
 F_n:D_n\longrightarrow D_n,\qquad
 P_n:D_{n+1}\longrightarrow D_n
\]
satisfy, for rational \(0\le q_n<1\), \(L_n\ge0\), and
\(\eta_n\ge0\),
\begin{align}
 d_n(F_nx,F_ny)&\le q_n d_n(x,y),
                                    \label{eq:level-contraction}\\
 d_n(P_nu,P_nv)&\le L_n d_{n+1}(u,v),
                                    \label{eq:restriction-lipschitz}\\
 d_n(F_nP_nz,P_nF_{n+1}z)&\le\eta_n.
                                    \label{eq:commutation-defect}
\end{align}
Let the actual finite computations return \(y_n\in D_n\) with
\[
 d_n(y_n,F_ny_n)\le r_n.             \label{eq:level-residual}\tag{R1}
\]
In particular, \(y_n\) need not be an exact fixed point. It can be a
rounded iterate with a residual certificate checked after computation.

\begin{theorem}[Residual and commutation bound]
\label{thm:commutation}
For the preceding finite data,
\[
 d_n(y_n,P_ny_{n+1})\le e_n,
 \qquad
 e_n=\frac{r_n+\eta_n+L_nr_{n+1}}{1-q_n}.
 \label{eq:adjacent-error}\tag{R2}
\]
\end{theorem}
\begin{proof}
Set \(a=d_n(y_n,P_ny_{n+1})\). Insert the intermediate points
\(F_ny_n\), \(F_nP_ny_{n+1}\), and \(P_nF_{n+1}y_{n+1}\).
The four successive distances are bounded respectively by
\(r_n\), \(q_na\), \(\eta_n\), and \(L_nr_{n+1}\). Thus
\[
 a\le r_n+q_na+\eta_n+L_nr_{n+1}.
\]
Cancellation gives \((1-q_n)a\le r_n+\eta_n+L_nr_{n+1}\).
The denominator is a supplied positive rational, so division is an
ordinary rational operation. Equivalently, keep this last inequality
as the denominator-free endpoint.
\end{proof}

The result is asymmetric in the two contraction factors: it uses only
the lower-level factor \(q_n\). Reversing the comparison would require
a suitable lift and its own properties. No such lift is assumed here.

\subsection{Finite telescoping through restrictions}

Define \(P_{n,n}\) to be the identity and
\(P_{n,m+1}=P_{n,m}\circ P_m\). This fixes the orientation of every
restriction and gives coherence by definition. Put
\[
 L_{n,n}=1,\qquad L_{n,j}=\prod_{t=n}^{j-1}L_t\quad(j>n).
\]
For the iPRA proof, fix both endpoints and two input vector codes.
Run the two descending restriction trajectories in parallel, with the
current level \(m-t\). Induct on the rational inequality
\[
 d_{m-t}(z_t,z'_t)\le
 \left(\prod_{j=m-t}^{m-1}L_j\right)d_m(z_0,z'_0),
\]
with the two domain-membership bits. The step is one instance of (\ref{eq:restriction-lipschitz}).
All endpoints and input vectors remain numerical parameters. Thus the
finite composition has the asserted Lipschitz bound without an induction
on a universally quantified statement about all vectors.

\begin{theorem}[Finite refinement estimate]
\label{thm:finite-refinement}
For \(m\ge n\),
\[
 d_n(y_n,P_{n,m}y_m)
 \le\sum_{j=n}^{m-1}L_{n,j}e_j.
 \label{eq:refinement-sum}\tag{R3}
\]
In particular, if every \(L_j\le1\), then
\[
 d_n(y_n,P_{n,m}y_m)\le\sum_{j=n}^{m-1}e_j.
\]
\end{theorem}
\begin{proof}
The case \(m=n\) has zero distance and empty sum. For the successor
step, insert \(P_{n,m}y_m\), apply the induction hypothesis to the
first distance, and use the Lipschitz bound of \(P_{n,m}\) and (R2)
for the second distance. This adds precisely \(L_{n,m}e_m\).
The induction is on a finite rational inequality. Its coded form is
quantifier-free.
\end{proof}

This proof does not take an infinite sum. A tail modulus means a PR
function \(T(n)\) with a separately proved finite assertion
\[
 \sum_{j=n}^{m-1}e_j\le T(n)\quad(m\ge n).
 \label{eq:finite-tail}\tag{R4}
\]
For example, \(e_j\le C2^{-j}\) gives \(T(n)=2C2^{-n}\) by the
finite geometric identity. When restrictions expand, the required
bound applies instead to the weighted sum in (R3). A bound on the
unweighted errors would then be insufficient.

\subsection{Producing named observations}

An observation index \(i\) specifies a minimal level \(b(i)\).
For \(n\ge b(i)\), suppose a PR observer \(O_{n,i}:D_n\to\Q\)
satisfies
\begin{align*}
 |O_{n,i}(x)-O_{n,i}(y)|&\le B_i d_n(x,y),\\
 O_{n,i}(P_nz)&=O_{n+1,i}(z),
\end{align*}
with a known rational \(B_i\ge0\). Exact observational
compatibility can be replaced by an explicit adjacent error, as
described below.

\begin{theorem}[A coordinate sampler from a refinement tower]
\label{thm:tower-sampler}
Assume nonexpansive restrictions and (R4). Suppose a PR schedule
\(\nu(i,k)\ge b(i)\), nondecreasing in \(k\), satisfies
\[
 B_iT(\nu(i,k))\le2^{-k-2}.
\]
Then
\[
 X(i,k)=O_{\nu(i,k),i}(y_{\nu(i,k)})
\]
is a PR family of regular rational names. More precisely,
\[
 |X(i,k)-X(i,l)|\le2^{-\min(k,l)-2}.
 \label{eq:tower-name}\tag{R5}
\]
\end{theorem}
\begin{proof}
If \(k\le l\), set \(n=\nu(i,k)\), \(m=\nu(i,l)\).
The compatibility equations, composed finitely many times, identify
\(O_{m,i}(y_m)\) with \(O_{n,i}(P_{n,m}y_m)\). Now apply the
Lipschitz bound of the observer, (R3), (R4), and the schedule inequality.
The case \(l\le k\) is symmetric. Natural order is decidable, so
the case split is intuitionistically valid.
\end{proof}

If instead
\(\abs{O_{n,i}(P_nz)-O_{n+1,i}(z)}\le\zeta_{i,n}\), the adjacent
observation difference is at most \(B_ie_n+\zeta_{i,n}\). A finite
tail bound for that sum gives exactly the same argument. Thus observers
may themselves be approximations, provided their errors are budgeted.

\subsection{How residual budgets are chosen}

The factor \((1-q_n)^{-1}\) is a numerical resource. For example,
suppose \(L_n\le1\), positive rationals \(a_n\) have a known
summable tail, and the input data satisfy
\[
 \eta_n\le\frac{1-q_n}{2}a_n,
 \qquad
 r_n+r_{n+1}\le\frac{1-q_n}{2}a_n.
\]
Then \(e_n\le a_n\). The first condition is a property of the
models and their restrictions; the second can be enforced by running
the finite solvers accurately enough. It is not possible to repair an
arbitrary model defect merely by performing more iterations of each
solver. More iterations reduce \(r_n\), not \(\eta_n\).

When gaps vary, one conservative local choice for a desired error
majorant \(a_n\) is to choose each \(r_n\), for \(n>1\), below
both \((1-q_n)a_n/4\) and \((1-q_{n-1})a_{n-1}/4\), and to
require \(\eta_n\le(1-q_n)a_n/2\). Endpoints are handled using
only the available neighboring condition. These are finite positive
rational budgets, and the approximate-action producer gives a bounded
precision and iteration schedule for meeting them whenever its
local certificates are available.

\subsection{Why this is a useful embedding interface}

The theorem permits one to work entirely with finite models, then
derive coherent observations without assuming an inverse-limit element
or choosing one exact solution at every level. It is applicable to
truncated coupled recurrences, finite sections of linear systems, and
certified spatial or spectral approximations when their actual
commutation defects can be bounded. These are application patterns,
not a claim that every PDE or every spectral scheme automatically
supplies those estimates. The next section supplies all the data for
one infinite nonlinear example.
